\section{Fixed-horizon limits and admissibility}\label{sec:admissibility}

In this section the physical horizon \(T\in(0,T_0]\) is fixed. The
approximation index tends to infinity without any rescaling of time or
space. This distinction matters: the cold densities in this section have
uniformly integrable limits, whereas their rescalings in the next section
will converge to measures in velocity.

\begin{proposition}\label{prop:admissible-limits}
Fix \(T\in(0,T_0]\). Either the zero-seed approximations of
Proposition~\ref{prop:finite-cap}, or any sequence of its cap shots on
\([0,T]\), has a subsequence whose total densities \(F_N=c_N+h_N\)
converge weakly in \(L^1([0,T]\times\Omega)\) to a nonnegative density
\(F\). The subsequence can be chosen so that
\(F_N(t)\rightharpoonup F(t)\) in \(L^1(\Omega)\) at every
\(t\in[0,T]\), with uniformly small polynomially weighted velocity
tails. The limit has the common initial datum \eqref{eq:datum}, and
\[
Q^\pm(F,F)\in L^1([0,T]\times\Omega),\qquad
F\in C([0,T];L^1(\Omega)).
\]
It satisfies every condition of Definition~\ref{def:admissible} on this
interval, including the inequalities at every time. Each such limit has
an admissible global continuation. The dormant upper bound, consisting
of the cold majorants of Proposition~\ref{prop:cold-control} and
\(t^2H\), passes to a zero-seed limit.
\end{proposition}

We give the compactness and entropy arguments in detail.

\subsection{Conservation and entropy before passage to the limit}

At a fixed approximation the bounds already proved imply
\(F_N\leq A_N H\) on the whole solution interval, with a finite constant
allowed to depend on \(N,b_N,T\). In particular there are all polynomial
velocity moments, and there is an \(L^p\) majorant for some
\(1<p<2/q\). The equation holds absolutely continuously on almost every
free path. Its cutoff gain and loss are integrable, including after
multiplication by any fixed polynomial in velocity.

The collision map preserves volume, kinetic energy, and
\(|v-v_*|\), so it also preserves the cutoff. For an integrable velocity
test \(\psi\) the collision symmetry gives
\[
\int_{\R^3}\psi(v)Q^{b_N}(F_N,F_N)(v)\dd v
=\frac12\int_{\R^3\times\R^3\times S^2}
 |a|\ind_{\{|v-v_*|<b_N\}}F_NF_{N,*}
 (\psi'+\psi_*'-\psi-\psi_*)\dd v\dd v_*\dd\omega.
\]
The right side vanishes for \(\psi=1,v_j,|v|^2\).
Integrating the path equation and using periodicity therefore proves
mass, momentum, and energy conservation at every time of the finite
approximation. The same calculation with a spatially dependent mass test
gives its local continuity equation.

To justify entropy at zeros, put \(E(v)=e^{-2|v|^2}\), take
\(0<\delta\leq1\), and set \(F_{N,\delta}=F_N+\delta E\). Define
\[
\mathcal A_\delta(f,v)
=M(v)+\int_0^f\log\frac{u+\delta E(v)}{M(v)}\dd u.
\]
Its derivative with respect to \(f\) is
\(\log(F_{N,\delta}/M)\). At the fixed approximation,
\begin{equation}\label{comp:log-bound}
\left|\log\frac{F_{N,\delta}}M\right|
\leq C_N\bigl(1+|\log\delta|+|v|^2
                   +\log^+|\bar v|^{-1}\bigr).
\end{equation}
Indeed \(F_{N,\delta}\geq\delta e^{-2|v|^2}\) gives the lower
logarithmic bound, and \(F_N\leq A_NH\) gives the upper bound.
Every factor on the right of \eqref{comp:log-bound} is integrable
against \(H\), apart from the harmless constant \(|\log\delta|\): in
the transverse plane this uses
\(\int_0^1 r^{1-q}|\log r|\dd r<\infty\).

Since \(E\leq C H\), the fixed-cutoff bilinear bounds imply
\[
|Q^{b_N}(F_{N,\delta},F_{N,\delta})
            -Q^{b_N}(F_N,F_N)|\leq C_{N,b_N}\delta H.
\]
Consequently replacing the latter operator by the former in its
integral against \(\log(F_{N,\delta}/M)\) costs at most
\(C_{N,b_N,T}\delta(1+|\log\delta|)\).
The path chain rule, followed by spatial integration, now gives
\begin{equation}\label{comp:regularized-entropy}
\int_\Omega\mathcal A_\delta(F_N(t),v)
 +\int_0^t\mathcal D_{b_N}(F_{N,\delta}(s))\dd s
\leq\int_\Omega\mathcal A_\delta(F_N(0),v)
       +C_{N,b_N,T}\delta(1+|\log\delta|).
\end{equation}
Here \(\mathcal D_b\) denotes the dissipation with the invariant
cutoff inserted. The collision identity used in this step is
\[
\int_\Omega Q^{b_N}(F_{N,\delta},F_{N,\delta})
       \log(F_{N,\delta}/M)
=-\mathcal D_{b_N}(F_{N,\delta}).
\]
It follows by averaging over interchange of the two particles and over
the elastic substitution. The \(\log M\) term cancels because it is a
linear combination of mass and energy. All integrals for this calculation
are finite at \(\delta>0\) by \eqref{comp:log-bound}.

As \(\delta\downarrow0\), the primitive converges to
\(f\log(f/M)-f+M\). Its integrals converge as well. For example its
lower bound is that limiting primitive, while its upper bound follows
by replacing \(\delta\) by one in the increasing integrand. The
negative part of \(f\log f\) is integrable under a finite second
moment, and the upper bound is integrable by \eqref{comp:log-bound}
without the \(|\log\delta|\) term. Fatou's lemma applied to the
nonnegative dissipation integrand in \eqref{comp:regularized-entropy}
therefore yields, for every \(t\in[0,T]\),
\begin{equation}\label{comp:finite-entropy}
\mathcal H_M(F_N(t))+\int_0^t\mathcal D_{b_N}(F_N(s))\dd s
\leq\mathcal H_M(F_N(0)).
\end{equation}

Lemma~\ref{lem:datum} gives convergence of the initial mass, energy,
and relative entropy to those of \eqref{eq:datum}. Its estimates include
the hot seed, whose mass and energy are \(O(s_N^4)\) and whose absolute
entropy is \(O(s_N^4(1+|\log s_N|))\), as well as the seed's overlap
with the cold density in the relative entropy. The convergence is
uniform for the allowed seed parameters \(0\leq\lambda\leq L_0\).

\subsection{Uniform integrability of densities and collision products}

For \(s=s_j\), let \(A_s\) be the annulus
\(\{y\in\R^2:s<|y|<e^\eta s\}\), viewed as a subset of \(\T^2\).
By Lemma~\ref{jet:geometry}, these annuli are pairwise disjoint and
contain the enlarged stripes of both signs at their respective scales.
Their areas satisfy \(|A_s|\leq Cs^2\). Proposition~\ref{prop:cold-control},
the cap, and the dormant bound imply, for every fixed \(k\geq0\),
\begin{equation}\label{comp:moment-envelope}
R_{k,N}(t,x):=\int_{\R^3}\langle v\rangle^kF_N(t,x,v)\dd v
\leq C_k\left(1+\sum_{j\leq N}
                  \frac{\ind_{A_{s_j}}(y)}{t+s_j}\right).
\end{equation}
The bounded term includes the hot density and escaped tails. The cold
remainder has density coefficient \(O(d_s/s^2)=O(s^8)\), and is
also covered by this estimate. All constants are independent of \(N\).

The square of the envelope is integrable uniformly in \(N\), including
at the initial surface. In fact the contribution of one annulus is
\begin{equation}\label{comp:annular-square}
\int_0^T\frac{Cs^2}{(t+s)^2}\dd t
=Cs^2\left(\frac1s-\frac1{T+s}\right)\leq Cs.
\end{equation}
Since the scales form a geometric sequence,
\(\sum_{s_j<\varepsilon}s_j\leq C_\eta\varepsilon\). If
\(A_{<\varepsilon}=\bigcup_{s_j<\varepsilon}A_{s_j}\), disjointness
of the annuli thus gives
\begin{equation}\label{comp:small-annuli}
\sup_N\int_0^T\int_{A_{<\varepsilon}\times\T}
                 R_{k,N}(t,x)^2\dd x\dd t\longrightarrow0
\qquad(\varepsilon\downarrow0).
\end{equation}
The contribution of the bounded term is at most
\(C T|A_{<\varepsilon}|\).

Outside \(A_{<\varepsilon}\), the main color has one of finitely many
scales \(s_j\geq\varepsilon\). Those
velocity widths are positive and fixed. Their Gaussian bounds, with
axial centers in a fixed bounded interval, give a function
\(\Psi_\varepsilon(v)\), independent of \(N,t,x\), such that
\begin{equation}\label{comp:fixed-majorant}
F_N(t,x,v)\leq\Psi_\varepsilon(v)
\quad\hbox{if }y\notin A_{<\varepsilon},\qquad
\int\langle v\rangle^k\Psi_\varepsilon(v)\dd v<\infty
\quad(k\geq0).
\end{equation}
One may take a finite sum of broadened Gaussians about the bounded axial
interval, together with a multiple of \(H\) for the hot component and
escaped tails. There is no need for this majorant to be bounded on the
velocity axis.

Equations \eqref{comp:small-annuli}--\eqref{comp:fixed-majorant} prove
uniform integrability of
\[
\langle v\rangle^k\langle w\rangle^k
 F_N(t,x,v)F_N(t,x,w)
\]
in the full spacetime and two-velocity variables. Indeed its integral
over the removed annuli is at most the left side of
\eqref{comp:small-annuli}; elsewhere it has the fixed integrable
majorant
\(\langle v\rangle^k\Psi_\varepsilon(v)
  \langle w\rangle^k\Psi_\varepsilon(w)\).
This also proves uniformly small velocity tails. Taking first powers
instead proves uniform integrability of the one-particle densities in
spacetime. At individual times the estimate is
\begin{equation}\label{comp:fixed-time-small}
\sup_{N,\,0\leq t\leq T}
 \int_{A_{<\varepsilon}\times\T}R_{k,N}(t,x)\dd x
\leq C_k\sum_{s_j<C\varepsilon}\frac{s_j^2}{s_j}
       +C_k|A_{<\varepsilon}|\longrightarrow0.
\end{equation}
Together with \eqref{comp:fixed-majorant}, this proves weighted
one-particle uniform integrability and velocity tightness uniformly at
all individual times.

The gains and losses enjoy the same spacetime conclusion. To see this
directly, collision symmetry and energy conservation imply
\[
\int\langle v\rangle^k
 Q_N^\pm(t,x,v)\dd v\leq C_k R_{k+1,N}(t,x)^2,
\qquad Q_N^\pm=Q^{\pm,b_N}(F_N,F_N).
\]
Their weighted integrals on the small annuli therefore tend uniformly
to zero. On the complement, positivity and \eqref{comp:fixed-majorant}
give the fixed bounds
\[
Q_N^+(t,x,v)\leq P(\Psi_\varepsilon,\Psi_\varepsilon)(v),\qquad
Q_N^-(t,x,v)\leq\Psi_\varepsilon(v)\nu_{\Psi_\varepsilon}(v).
\]
Both right sides are integrable with every fixed polynomial velocity
weight: integrate first in the output velocity and use the preceding
collision-moment inequality for the fixed function
\(\Psi_\varepsilon\). Thus \(\{Q_N^\pm\}\) is uniformly integrable
in spacetime, with uniformly small weighted velocity tails. In
particular its integrals over time intervals of length tending to zero
vanish uniformly in \(N\). This last assertion follows by first
discarding the small annuli and then integrating the fixed majorants
over the time interval.

\subsection{An elementary averaging argument}

We use a uniformly integrable version of velocity averaging
\cite[Proposition~3]{GolseLionsPerthameSentis1988}. The following proof gives
the precise local statement needed here.

\begin{lemma}\label{comp:averaging}
Suppose that \(f_n\) and \(g_n\) are locally uniformly integrable in
\(L^1\) on an open time-space-velocity region, and
\((\partial_t+v\cdot\nabla_x)f_n=g_n\) there in distributions.
Then for every smooth compactly supported velocity test \(\psi\), the
averages \(\int\psi(v)f_n(t,x,v)\dd v\) are relatively compact in
local \(L^1_{t,x}\) on regions where the hypotheses hold for velocities
in a neighborhood of \(\supp\psi\). The assertion applies both on a
spatial torus and locally in \(\R^3\).
\end{lemma}

\begin{proof}
Choose smooth interior cutoffs in time, space, and velocity; the spatial
cutoff is unnecessary on a torus. Write the localized sequence as
\(q_n\), extend it by zero outside the localization interval, and put
\(r_n=(1+\partial_t+v\cdot\nabla_x)q_n\). The equation shows that
\(r_n\) is uniformly integrable in \(L^1\), with support in a fixed
bounded set of time and velocity and, in the Euclidean case, space.
The derivative of a spatial cutoff is harmless because velocity has
also been localized. The exact identity
\[
q_n(t,x,v)=\int_0^\infty e^{-s}r_n(t-s,x-sv,v)\dd s
\]
follows by integrating the localized equation from a time preceding its
support. Denote this damped inverse by \(R\). It is an \(L^1\)
contraction, since free transport preserves Lebesgue measure.

Truncate \(r_n\) at height \(L\), retaining its sign, to obtain
\(r_n^{(L)}\). Uniform integrability gives
\[
\sup_n\|r_n-r_n^{(L)}\|_1\longrightarrow0,
\qquad
\sup_n\|r_n^{(L)}\|_2<\infty
\quad\hbox{for fixed }L.
\]
The averages of \(R(r_n-r_n^{(L)})\) are therefore uniformly small
in \(L^1\). It remains to handle a bounded \(L^2\) family \(r_n^{(L)}\).

Take Fourier transform in \((t,x)\), using Fourier series in the
periodic coordinates. For \(u_n=Rr_n^{(L)}\), Cauchy--Schwarz gives
\[
\left|\int\psi(v)\widehat u_n(\tau,k,v)\dd v\right|^2
\leq I_\psi(\tau,k)
       \int|\widehat r_n^{(L)}(\tau,k,v)|^2\dd v,
\qquad
I_\psi(\tau,k)=\int
 \frac{|\psi(v)|^2}{1+(\tau+k\cdot v)^2}\dd v.
\]
If \(\supp\psi\subset\{|v|\leq R_0\}\), integration in the
direction of \(k\) gives \(I_\psi\leq C_\psi/|k|\) when \(k\ne0\).
If \(|\tau|>2R_0|k|\), the denominator instead gives
\(I_\psi\leq C_\psi/(1+|\tau|^2)\). It follows that
\(I_\psi(\tau,k)\to0\) uniformly as \(|\tau|+|k|\to\infty\).
The averages consequently have uniformly small high-frequency
\(L^2\) tails. Convolution with a fixed smooth approximate identity
approximates them uniformly in \(L^2\). On each compact time-space
set the convolutions are bounded and equicontinuous, by Cauchy--Schwarz
applied to translates of the convolution kernel and their derivatives.
They are therefore relatively compact there. This proves local
\(L^2\), hence local \(L^1\), compactness for fixed \(L\).
The uniform \(L^1\) error as \(L\to\infty\) proves the assertion.
\end{proof}

\subsection{Identification of the nonlinear terms}

The uniform integrability just proved supplies a subsequence with
\(F_N\rightharpoonup F\) in spacetime \(L^1\); velocity tightness
allows the whole velocity space. Since
\(\D F_N=Q_N^+-Q_N^-\), Lemma~\ref{comp:averaging} applies.
After a diagonal extraction for a countable dense family of compact
velocity tests, all such velocity averages converge strongly in local
\(L^1_{t,x}\) to the corresponding averages of \(F\). Extension from
that family to every smooth compact velocity test uses the uniform
mass bound.

Set \(W_N(t,x,v,w)=F_N(t,x,v)F_N(t,x,w)\). The full pair-product
uniform integrability supplies a weak \(L^1\) limit \(W\), after
another extraction. For smooth compact velocity tests \(\psi,\chi\),
\[
\int\psi(v)\chi(w)W_N\dd v\dd w
=\left(\int\psi F_N\dd v\right)
 \left(\int\chi F_N\dd w\right).
\]
Both factors converge in local \(L^1\), so their product converges in
measure. The products are uniformly integrable in \((t,x)\), by
\eqref{comp:annular-square} and the fixed majorants after removal of
small annuli. Their convergence is consequently in \(L^1\): to verify
this last implication, discard a set of arbitrarily small measure on
which convergence is not uniform, and use uniform integrability on the
discarded set. We obtain
\[
\int\psi(v)\chi(w)W\dd v\dd w
=\left(\int\psi F\dd v\right)
 \left(\int\chi F\dd w\right).
\]
Finite sums of separated velocity tests are uniformly dense among
continuous functions on compact two-velocity sets. Testing also in
\((t,x)\) identifies
\begin{equation}\label{comp:pair-limit}
W_N\rightharpoonup F(t,x,v)F(t,x,w)
\quad\hbox{in }L^1([0,T]\times\T^3\times\R^3\times\R^3).
\end{equation}
The same conclusion holds with any fixed polynomial weight for which
it is needed. Interior time compactness suffices: the uniform
integrability already proved removes intervals near zero and \(T\).

Adjoining the finite angular measure preserves weak convergence. The
elastic map on \((v,w,\omega)\) is measure preserving and involutive,
so \eqref{comp:pair-limit} also identifies the primed pair product.
On a bounded two-velocity set the factor \(|a|\) is bounded, and the
cutoff is eventually absent. The weighted tail estimates allow these
sets to exhaust all velocities. Since \(Q_N^\pm\) themselves are
uniformly integrable, testing their weak limits and using these
identifications proves
\begin{equation}\label{comp:collision-limit}
Q_N^\pm\rightharpoonup Q^\pm(F,F)
\quad\hbox{in }L^1([0,T]\times\Omega).
\end{equation}
Passing the equation, including its initial boundary term, is now
legitimate. The initial densities converge in \(L^1\), and the limit
satisfies the unrenormalized distributional equation with datum \(F_0\).

\subsection{Every-time traces, entropy, and the chain rule}

We first specify representatives at every time. For a smooth periodic
test \(\varphi(x,v)\) compactly supported in velocity, the finite path
equations give
\[
|\langle F_N(t)-F_N(s),\varphi\rangle|
\leq\int_s^t\int_\Omega
 \bigl(F_N|v\cdot\nabla_x\varphi|
       +(Q_N^++Q_N^-)|\varphi|\bigr)\dd x\dd v\dd u.
\]
The first term is bounded by \(C_\varphi|t-s|\); the second has a
uniform modulus tending to zero with \(|t-s|\), by the time-integral
conclusion above. For a countable uniformly dense collection of such
tests, extract uniform convergence of their pairings on \([0,T]\).
At each individual time \eqref{comp:fixed-time-small} and
\eqref{comp:fixed-majorant} make \(\{F_N(t)\}\) relatively weakly
compact in \(L^1\). Every cluster point has the prescribed smooth
pairings, hence is the same density. Thus this single subsequence
converges weakly in \(L^1\) at every time. It defines a weakly
continuous representative against smooth tests, agreeing with the
spacetime limit. The weighted fixed-time tightness also passes total
mass, momentum, and energy from the finite approximations. In
particular energy is conserved on this early interval, although only
its inequality will be needed.

Let
\(\Phi(A,B)=(A-B)\log(A/B)\) with its nonnegative extended values.
It is jointly convex and lower semicontinuous: it is the sum of the
two convex perspectives \(A\log(A/B)\) and \(B\log(B/A)\).
On every bounded collision set, \eqref{comp:pair-limit} and its primed
version therefore give lower semicontinuity of the integral of
\(|a|\Phi\). One can absorb the weight into the arguments because
\(\Phi(cA,cB)=c\Phi(A,B)\) for \(c\geq0\). The lower
semicontinuity follows, for example, by expressing a nonnegative convex
lower-semicontinuous function as the supremum of its affine supporting
functions and using weak convergence on a finite partition. Exhausting
bounded velocity sets, and using that the cutoff is eventually absent
on each of them, gives for every \(t\in[0,T]\)
\[
\int_0^t\mathcal D(F(s))\dd s
\leq\liminf_N\int_0^t\mathcal D_{b_N}(F_N(s))\dd s.
\]
The same convex lower-semicontinuity argument applies to
\(\mathcal H_M(F_N(t))\) under weak \(L^1\) convergence at that
time. Combining the two nonnegative terms, using
\eqref{comp:finite-entropy} and initial entropy convergence, proves
\begin{equation}\label{comp:limit-entropy}
\mathcal H_M(F(t))+\int_0^t\mathcal D(F(s))\dd s
\leq\mathcal H_M(F_0)\qquad(0\leq t\leq T).
\end{equation}
There is no time-dependent subsequence in this conclusion.

In fact the time representative is strongly continuous. In free
coordinates the distributional equation and its initial trace read
\[
\partial_t\bigl[F(t,x+tv,v)\bigr]=Q(F,F)(t,x+tv,v).
\]
The right side is integrable in \((t,x,v)\). Subtracting its Bochner
integral leaves a distribution with zero time derivative, whose initial
value is \(F_0\). Therefore
\begin{equation}\label{comp:path-formula}
F(t,x+tv,v)=F_0(x,v)
 +\int_0^t Q(F,F)(s,x+sv,v)\dd s
\end{equation}
as an \(L^1_{x,v}\) identity, initially for almost every time and then
for every time by the representative just constructed. The right side
is absolutely continuous in \(L^1\). Free transport is a strongly
continuous group on \(L^1(\T^3\times\R^3)\), as follows by
approximating an \(L^1\) function by continuous functions with compact
velocity support. Thus \(F\in C([0,T];L^1)\). This proves in
particular continuity against every bounded measurable test.

Equation \eqref{comp:path-formula} also holds on almost every individual
free path. Since \(M(v)>0\) is constant on such a path, the ordinary
absolutely continuous chain rule gives, for every \(\beta\) in
Definition~\ref{def:admissible},
\[
\D\bigl[M\beta(F/M)\bigr]=\beta'(F/M)Q(F,F).
\]
The source is integrable because \(|\beta'|\leq C_\beta\) and
\(Q\in L^1\). Moreover
\[
|M\beta(f/M)-M\beta(g/M)|\leq C_\beta|f-g|,
\qquad |M\beta(f/M)|\leq C_\beta f.
\]
Hence the renormalization is strongly continuous at the initial and
terminal traces. Integrating its path equation proves exactly the
renormalized identity and trace in Definition~\ref{def:admissible}.
The separate integrability requirement follows from
\(0\leq Q^\pm/(1+F/M)\leq Q^\pm\).

Integrating the unrenormalized equation in velocity is justified by
absolute collision integrability. The collision integral has zero mass,
and the flux is integrable, so it gives the required local continuity
equation with its initial density. Momentum has already passed at every
time. Finally, the negative part of \(F\log F\) is bounded uniformly
by mass and energy. For example split \(\{F<1\}\) at
\(F=e^{-1-|v|^2}\): above that threshold use
\(-F\log F\leq(1+|v|^2)F\), and below it use a fixed multiple of
\(e^{-(1+|v|^2)/2}\). Equation \eqref{comp:limit-entropy} then bounds
the positive part as well. This proves the uniform absolute logarithmic
integral required in the definition.

\subsection{Continuation and inherited bounds}

At time \(T\), the density just constructed has finite mass, energy,
and absolute logarithmic integral. We use the global existence theorem
of DiPerna and Lions \cite{DiPernaLions1989,DiPernaLions1991}, in the
precise periodic formulation of
\cite[Section~4.1, Theorem~4.1]{LevermoreMasmoudi2010}. We verify its
interface with Definition~\ref{def:admissible} before concatenating.

For our kernel \(b_{\mathrm{phys}}(\omega,v-w)=|(v-w)\cdot\omega|\),
local integrability holds and, for every compact \(K\subset\R^3\),
\[
\frac1{1+|v|^2}\int_{K\times S^2}
 b_{\mathrm{phys}}(\omega,v-w)\dd\omega\dd w
=\frac{2\pi}{1+|v|^2}\int_K|v-w|\dd w\longrightarrow0
\quad(|v|\to\infty).
\]
This is the hypothesis of that theorem. To match its collision and time
normalizations, put
\[
C_b=\int_{\R^3\times\R^3\times S^2}
 b_{\mathrm{phys}}(\omega,v-w)M(v)M(w)\dd\omega\dd v\dd w,
\qquad \epsilon=C_b^{-1}.
\]
The constant \(C_b\) is finite and positive. Use the source kernel
\(b_{\mathrm{phys}}/C_b\) and its equation
\(\epsilon\partial_sG+v\cdot\nabla_xG
=\epsilon^{-1}Q_{b_{\mathrm{phys}}/C_b}(G,G)\), with initial
density \(F(T)/M\). Here the relative-density operator satisfies
\(M Q_{b_{\mathrm{phys}}/C_b}(G,G)=Q(MG,MG)/C_b\).
Then \(\widetilde F(t,x,v)=M(v)G(\epsilon t,x,v)\) has precisely
our collision normalization. The same time substitution converts the
source dissipation coefficient \(\epsilon^{-2}\) into the coefficient
one of \(\mathcal D\).

The cited theorem allows arbitrary finite relative-entropy initial
data. It gives weak continuity of \(G\) in \(L^1(M\dd x\dd v)\),
which is exactly weak continuity of \(\widetilde F\) in
\(L^1(\dd x\dd v)\), and the renormalized identities with the initial
trace for continuously differentiable functions satisfying
\(|\beta'(z)|\le C/\sqrt{1+z}\). Our renormalizations are included.
Its local mass identity is the one required here. Integrating its
momentum identity over the torus removes the spatial divergence of
the momentum defect, giving total momentum conservation. This extends
to every time by weak continuity and the uniform second-moment bound:
the first-moment tail outside \(|v|\le R\) is at most \(R^{-1}\)
times the second moment. Dropping the nonnegative defect contributions
in the theorem's energy and relative-entropy inequalities gives the
two inequalities in Definition~\ref{def:admissible} at every time.
The negative-entropy estimate above then supplies the uniform absolute
logarithmic integral.

We also need separate global integrability of the renormalized gain and
loss, rather than only their difference. For any density \(f\ge0\),
write \(g=f/M\). Since \(f/(1+g)\le M\),
\[
\int_\Omega\frac{Q^-(f,f)}{1+g}
\le\int_\Omega M\nu_f
\le C\int_\Omega(1+|v|)f.
\]
For \(A,B\ge0\), with the same extended-value convention as for the
dissipation,
\[
A\le2B+\frac2{\log2}(A-B)\log(A/B).
\]
Indeed the first term suffices when \(A\le2B\); otherwise
\(A-B>A/2\) and \(\log(A/B)>\log2\). Apply this with
\(A=f'f_*'\), \(B=ff_*\), divide by \(1+g(v)\), and integrate
with the collision kernel. Nonnegativity of the dissipation gives
\[
\int_\Omega\frac{Q^+(f,f)}{1+g}
\le2\int_\Omega\frac{Q^-(f,f)}{1+g}
  +\frac8{\log2}\mathcal D(f).
\]
The moment and entropy-dissipation bounds for \(\widetilde F\) make
both terms integrable over every finite time interval. Thus the
continuation satisfies all of Definition~\ref{def:admissible}.

Concatenate \(\widetilde F\) with \(F\) at time \(T\).
The two weak \(L^1\) traces agree, so the concatenation is
continuous against every bounded measurable test. Its renormalized
identity follows by adding the two identities: the early terminal term
\(M\beta(F(T)/M)\) cancels the continuation's prescribed initial
term. The Lipschitz estimate above supplies that early terminal trace.
The mass continuity identities concatenate in the same way.

For \(t\geq T\), add the continuation inequality
\[
\mathcal H_M(\widetilde F(t-T))
 +\int_0^{t-T}\mathcal D(\widetilde F(s))\dd s
\leq\mathcal H_M(F(T))
\]
to \eqref{comp:limit-entropy} at \(T\). This gives the global entropy
inequality, with no entropy equality at the splice assumed. Energy
inequalities and momentum conservation concatenate using the common
terminal moments. Collision integrability in the renormalized sense
holds on both pieces, hence on every bounded interval of the
concatenation. The uniform moment and absolute entropy bounds on such
an interval follow from those of the two pieces.

For a dormant sequence, sum the cold upper bounds of
Proposition~\ref{prop:cold-control} over all labels and add \(t^2H\).
This is a common, locally integrable pointwise upper bound for the
finite total densities. Testing against nonnegative functions and
passing to the weak limit preserves it. All assertions of
Proposition~\ref{prop:admissible-limits} are now proved.
